\documentclass[12pt]{article}
\usepackage[utf8]{inputenc}
\usepackage[T1]{fontenc}
\usepackage{hyperref}
\title{Venetic Language Resources for K-12 Curriculum Development}
\author{Compiled by Claude Sonnet 5.5 (claude-sonnet-5-5) for Andr\'as Kornai}
\date{}
\begin{document}
\maketitle

\section{Introduction}
Venetic (ISO 639-3 xve; Glottocode vene1257) is an extinct Indo-European language of ancient northeastern Italy. It is known only from inscriptions, mostly from Este and Padova, numbering several hundred. Counts vary by source and by how fragments and Veneto-Latin texts are classed. There are no native speakers and no native-language educational materials. Any curriculum would be a heritage and epigraphy curriculum, not one taught in the language.

\section{Primary Educational Resources}
No primers, textbooks or school programmes for Venetic turned up in the searches. The scholarly pages below could be adapted by teachers.
\begin{itemize}
\item Developmental Literacy Case-Study 1: North East Italy (UCL Ancient Literacy project): \url{https://www.ucl.ac.uk/ancient-literacy/project-case1.htm}
\item What Did Palaeo-European Peoples Write? Local Languages of the Western Mediterranean (Antigone), an accessible overview: \url{https://antigonejournal.com/2023/05/western-mediterranean-ancient-languages/}
\end{itemize}

\section{Linguistic and Reference Materials}
\begin{itemize}
\item Marinetti, ``Venetico'' (2020, \emph{Palaeohispanica}), an overview of the corpus: \url{https://iris.unive.it/retrieve/e4239ddd-ff7b-7180-e053-3705fe0a3322/Marinetti%20Venetico%202020%20Paleohispanica.pdf}
\item PRIN Italia antica project page, University of Florence: \url{https://www.prin-italia-antica.unifi.it/p148.html}
\item Wikipedia, Venetic language: \url{https://wikipedia.com/wiki/Venetic}
\end{itemize}
Lejeune's \emph{Manuel de la langue v\'en\`ete} (1974) is the standard reference. The search results mention it, but I did not find a URL for it.

\section{Religious and Community Resources}
No living religious or community organisations tied to Venetic turned up in the searches.

\section{Digital and Audio Resources}
There are no audio resources, since the pronunciation is reconstructed. For the digital corpus effort, see:
\begin{itemize}
\item Katherine McDonald, Building a Corpus of Ancient Venetic: \url{https://katherinemcdonald.net/2015/08/25/building-a-corpus-of-ancient-venetic/}
\end{itemize}

\section{Curriculum Development Context}
Venetic is a language known only from inscriptions, so a K-12 programme could cover the script, selected inscriptions and the history of the Veneti, but not language immersion. The available material is scholarly and mostly aimed at adults. Teachers would need to adapt it.
\end{document}